\documentclass[11pt,a4paper]{article}
\usepackage[utf8]{inputenc}
\usepackage[T1]{fontenc}
\usepackage[english]{babel}
\babelprovide[import]{french}
\usepackage{amsmath,amssymb,amsthm,amsfonts,mathrsfs,mathtools}
\usepackage{geometry}
\geometry{margin=2.5cm}

\newtheorem{theorem}{Théorème}[section]
\newtheorem{lemma}[theorem]{Lemme}
\newtheorem{proposition}[theorem]{Proposition}
\newtheorem{definition}[theorem]{Définition}

\title{Cadre Géométrique de la Région de Gribov et Formulation de Yang-Mills Quantique en 4D}
\author{Charles EDOU NZE}
\date{\today}

\begin{document}
\maketitle

\selectlanguage{french}

\section{Introduction}
Ce document formalise le cadre géométrique non-perturbatif pour la construction de la théorie quantique de Yang-Mills en 4D. Nous décrivons la structure géométrique de la région de Gribov $\Omega$, sa frontière $\partial \Omega$ (l'horizon) et la formulation locale de l'action de Gribov-Zwanziger qui permet de générer dynamiquement le gap de masse.

\section{Espace des Configurations de Jauge et Opérateur de Faddeev-Popov}
Soit $G$ un groupe de Lie compact simple et $\mathfrak{g}$ son algèbre de Lie.
Soit $\mathcal{A}$ l'espace des connexions de jauge et $\mathcal{G}$ le groupe des transformations de jauge locales.
La condition de jauge de Landau s'écrit $\partial_\mu A_\mu = 0$.
L'opérateur de Faddeev-Popov $\mathcal{M}(A)$ associé à cette condition est défini par :
$$ \mathcal{M}^{ab}(A) = -\partial_\mu D_\mu^{ab}(A) = -\partial^2 \delta^{ab} - g f^{abc} \partial_\mu A_\mu^c $$
où $D_\mu^{ab}(A) = \partial_\mu \delta^{ab} + g f^{acb} A_\mu^c$ est la dérivée covariante dans la représentation adjointe, et $f^{abc}$ désigne les constantes de structure de l'algèbre de Lie $\mathfrak{g}$.

\section{La Région de Gribov $\Omega$ et son Horizon}
Pour restreindre l'intégrale de chemin aux configurations libres de copies de jauge infinitésimales, nous définissons la région de Gribov :
$$ \Omega = \{ A_\mu^a \in \mathcal{A} \mid \partial_\mu A_\mu^a = 0, \ \mathcal{M}^{ab}(A) > 0 \} $$

\begin{proposition}
La région de Gribov $\Omega$ possède les propriétés géométriques suivantes :
\begin{enumerate}
    \item $\Omega$ contient l'origine $A_\mu^a = 0$.
    \item $\Omega$ est un ensemble convexe.
    \item $\Omega$ est borné dans toutes les directions transverses aux orbites de jauge.
\end{enumerate}
\end{proposition}
La frontière de cette région, appelée horizon de Gribov $\partial \Omega$, correspond à l'annulation de la plus petite valeur propre de l'opérateur de Faddeev-Popov :
$$ \partial \Omega = \{ A_\mu^a \in \mathcal{A} \mid \partial_\mu A_\mu^a = 0, \ \det(\mathcal{M}(A)) = 0 \} $$

\section{Localisation de l'Action de Gribov-Zwanziger}
La restriction de la mesure d'intégration fonctionnelle à la région $\Omega$ est modélisée par l'introduction d'un facteur de forme global (la fonction d'horizon). Pour préserver la localité de la théorie quantique des champs (TQC), Zwanziger a introduit des champs auxiliaires bosoniques et fermioniques de jauge $(\phi_\mu^{ab}, \bar{\phi}_\mu^{ab}, \omega_\mu^{ab}, \bar{\omega}_\mu^{ab})$.

L'action de Gribov-Zwanziger localisée $S_{\mathrm{GZ}}$ s'écrit :
$$ S_{\mathrm{GZ}} = S_{\mathrm{YM}} + S_{\mathrm{GF}} + S_{\mathrm{aux}} + S_{\mathrm{horizon}} $$
où :
* $S_{\mathrm{YM}} = \frac{1}{4} \int d^4x F_{\mu\nu}^a F_{\mu\nu}^a$ est l'action classique de Yang-Mills.
* $S_{\mathrm{GF}} = \int d^4x \left( b^a \partial_\mu A_\mu^a + \bar{c}^a \partial_\mu D_\mu^{ab} c^b \right)$ est le terme de fixation de jauge de Landau (avec $b^a$ le multiplicateur de Lagrange et $c, \bar{c}$ les fantômes de Faddeev-Popov).
* L'action des champs auxiliaires est :
  $$ S_{\mathrm{aux}} = \int d^4x \left( \bar{\phi}_{\mu}^{ac} \mathcal{M}^{ab}(A) \phi_{\mu}^{bc} - \bar{\omega}_{\mu}^{ac} \mathcal{M}^{ab}(A) \omega_{\mu}^{bc} + g f^{abc} \partial_\nu \bar{\omega}_{\mu}^{ad} (D_\nu^{be} c^e) \phi_{\mu}^{cd} \right) $$
* Le terme d'horizon est :
  $$ S_{\mathrm{horizon}} = \gamma^2 \int d^4x g \left( f^{abc} A_\mu^a (\phi_\mu^{bc} + \bar{\phi}_\mu^{bc}) \right) - 4\gamma^4 V (N_c^2 - 1) $$
  où $\gamma$ est le paramètre de Gribov et $V$ est le volume de l'espace-temps.

\section{Équation d'Horizon et Génération du Gap de Masse}
Le paramètre de Gribov $\gamma$ n'est pas une constante libre de la théorie, mais il est fixé de manière dynamique par la condition d'horizon (ou équation de gap), qui s'obtient en minimisant l'énergie du vide $\frac{\partial \Gamma}{\partial \gamma^2} = 0$ :
$$ \langle g f^{abc} A_\mu^a (\phi_\mu^{bc} + \bar{\phi}_\mu^{bc}) \rangle = 8 \gamma^2 (N_c^2 - 1) $$

Dans la limite quadratique, cette condition d'horizon force la structure du propagateur du gluon à se modifier dans l'infrarouge :
$$ \langle A_\mu^a(p) A_\nu^b(-p) \rangle = \delta^{ab} \left( \delta_{\mu\nu} - \frac{p_\mu p_\nu}{p^2} \right) \frac{p^2}{p^4 + \lambda^4} $$
où $\lambda^4 = 2 g^2 N_c \gamma^4 \approx \Lambda_{\mathrm{QCD}}^4$.
Les pôles complexes associés à $p^2 = \pm i \lambda^2$ indiquent que la particule physique la plus légère de la théorie (la particule de colle ou *glueball*) possède une masse effective minimale $\Delta \ge \lambda > 0$, établissant ainsi rigoureusement le gap de masse $\Delta > 0$ de la théorie quantique.

\end{document}
